% ============================================================
%  sections/02_related_work.tex
%  2. Related Work (~0.75 page)
% ============================================================

\section{Related Work}
\label{sec:related}

\textbf{RAG evaluation frameworks.}
RAGAS introduced reference-free metrics for the retriever and the
generator separately, including faithfulness, answer relevance, and
context relevance, and validated them against human
judgements~\cite{es2023ragas}. It measures text quality in isolation and
does not account for latency or monetary cost. ARES adds lightweight
trained judges and evaluates under domain shift, improving on pure
prompting-based judges but still centering on answer quality~\cite{saadfalcon2023ares}.
eRAG evaluates retrieval quality in a way that correlates better with
downstream RAG performance than raw relevance labels, reinforcing that
retrieval and generation should be evaluated jointly but
separately~\cite{salemi2024erag}. RAGVUE decomposes evaluation into
claim-level and diagnostic components, making it possible to attribute
failures to retrieval, grounding, or reasoning~\cite{murugaraj2025ragvue}.
RARE stresses perturbation robustness and time-sensitive corpora,
showing that RAG systems can degrade under realistic distribution
shifts~\cite{zeng2025rare}. RGB reframes evaluation as a stress test of
four capabilities: noise robustness, negative rejection, information
integration, and counterfactual robustness, and reports rejection rate
alongside accuracy~\cite{chen2023rgb}. These frameworks score a
pipeline; most still treat cost and latency as secondary rather than
first-class axes.

\textbf{Corrective and self-correcting retrieval.}
CRAG adds a lightweight retrieval evaluator that scores each retrieved
document and triggers one of three actions---Correct, Incorrect, or
Ambiguous---then refines the context with a decompose-then-recompose step~\cite{yan2024crag}.
It improves robustness when retrieval degrades, but published overhead
numbers exclude retrieval and preprocessing, and the method assumes an
external evaluator checkpoint plus an optional web-search
fallback~\cite{yan2024crag}. The version evaluated here uses a local
proxy evaluator and no live search (Section~\ref{sec:design}); it is
therefore a derivative corrective design, not a full CRAG reproduction.

\textbf{Cost, latency, and component analysis.}
EvaRAG builds a large controlled configuration space over indexing,
similarity, retrieval, and reranking, and logs semantic quality together
with mean and P95 latency, token usage, and monetary cost per
query~\cite{elkiran2026empirical}. Rasheed et al.\ apply factorial ANOVA
and effect-size analysis to the same family of RAG components and show
that retriever and similarity choices dominate quality variance while
index choice mainly trades latency~\cite{rasheed2025empirical}. These
works establish the scorecard this paper adopts, but they study large
hosted models and treat retrieval as an experimental axis. Fixing
retrieval and moving to small open models changes the question from
``which pipeline component matters?'' to ``which deployable model should
I pick?''.

\textbf{Faithfulness in small models.}
Mala et al.\ benchmark faithfulness hallucinations for 2--8B open models
across sparse, dense, hybrid, and reranked retrieval, and show that
stronger retrieval can raise raw accuracy while also raising
hallucination rate; they introduce adjusted accuracy, which scores
precision over attempted answers only~\cite{mala2026faithfulness}.
Their study motivates our main faithfulness metric and our separation of
rejection rate from accuracy. Importantly, faithfulness here means
support by retrieved evidence, not factual truth; the two can diverge~\cite{murugaraj2025ragvue}.

Table~\ref{tab:positioning} summarises how existing work differs from
this benchmark.

% ============================================================
%  tables/tab_positioning.tex
%  Positioning table for Section 2
% ============================================================

\begin{table}[h]
\centering
\caption{Positioning of this benchmark relative to prior work.}
\label{tab:positioning}
\begin{tabular}{llllp{2.6cm}}
\toprule
\textbf{Work} & \textbf{Models} & \textbf{Retrieval axis} & \textbf{Core metrics} & \textbf{Missing} \\
\midrule
RAGAS~\cite{es2023ragas} & Hosted LLMs & Fixed (developer-chosen) & Faithfulness, relevance & Latency, cost \\
ARES~\cite{saadfalcon2023ares} & Hosted/SLM & Fixed & Trained judges, domain shift & Latency, cost \\
eRAG~\cite{salemi2024erag} & Hosted/SLM & Fixed & Retrieval quality correlation & Cost, SLMs \\
RAGVUE~\cite{murugaraj2025ragvue} & Hosted/SLM & Fixed & Claim-level diagnostics & Latency, cost \\
RARE~\cite{zeng2025rare} & Hosted/SLM & Fixed & Perturbation robustness & Cost \\
RGB~\cite{chen2023rgb} & LLMs (EN/ZH) & Fixed per setup & Accuracy, rejection rate & Cost, small models \\
CRAG~\cite{yan2024crag} & Off-the-shelf LLMs & Fixed per setup & Accuracy gains & Latency, cost \\
Mala et al.~\cite{mala2026faithfulness} & 2--8B open models & Varied (sparse/dense/hybrid) & Accuracy, hallucination, rejection & Latency, cost \\
EvaRAG~\cite{elkiran2026empirical} & Hosted models & Varied across 162 configs & Correctness, cost, latency & Faithfulness on SLMs \\
Rasheed et al.~\cite{rasheed2025empirical} & Hosted models & Varied & Correctness, cost, latency & SLM faithfulness \\
This work & Qwen2.5-1.5B / Qwen3-4B / LLaMA-3.1-8B & \textbf{Fixed (bm25s $\to$ MiniLM)} & \textbf{Faithfulness + latency + cost + abstention} & --- \\
\bottomrule
\end{tabular}
\end{table}

